% Added in the July 2026 review pass (annotation [155]): step-by-step derivations of the
% geodetic-to-ECEF transform (eq:ecef) and of the ECEF-to-ENU rotation used in
% Sec.~\ref{sec:enu}, plus the tangent-plane error estimate quoted there. Entirely new
% content, hence wrapped in \rev{} as one block.
\chapter{Geodetic transforms: from \texorpdfstring{$(\varphi,\lambda,h)$}{(phi, lambda, h)} to ECEF to ENU}
\label{app:geodesy}

\begin{revblock}
This appendix derives the two formulas that Sec.~\ref{sec:enu} uses as given: the
geodetic-to-ECEF transform, Eq.~\eqref{eq:ecef}, and the ECEF-to-ENU rotation matrix. It
closes with the size of the tangent-plane error quoted there. Notation as in the body:
WGS84 semi-major axis $a$, flattening $f$, $b=a(1-f)$ the semi-minor axis,
$e^2=f(2-f)=1-b^2/a^2$ the first eccentricity squared~\cite{nga_wgs84}.

\section{Geodetic coordinates to ECEF}
\label{app:geodesy-ecef}

The ECEF frame is Cartesian, centred at the Earth's centre $O$: the $Z$ axis along the
polar axis, the $X$ axis through the intersection of the equator and the prime meridian,
$Y$ completing the right-handed triad. By symmetry, a point at longitude $\lambda$ lies in
the meridian plane obtained by rotating the $X$--$Z$ plane by $\lambda$; writing $p$ for
the distance from the polar axis, its ECEF coordinates are
\begin{equation}
\label{eq:geo-meridian}
X=p\cos\lambda,\qquad Y=p\sin\lambda,\qquad Z=Z .
\end{equation}
The problem therefore reduces to the meridian ellipse: find $(p,Z)$ of a point at geodetic
latitude $\varphi$ and height $h$.

\paragraph{The surface point.}
The meridian section of the ellipsoid is the ellipse
\begin{equation}
\label{eq:geo-ellipse}
\frac{p^2}{a^2}+\frac{Z^2}{b^2}=1 .
\end{equation}
The geodetic latitude is defined through the surface \emph{normal}
(Figure~\ref{fig:ellipsoid}). The gradient of the left-hand side of
\eqref{eq:geo-ellipse} is proportional to $(p/a^2,\;Z/b^2)$, so the outward normal at a
surface point $(p_s,Z_s)$ makes an angle $\varphi$ with the equatorial plane given by
\begin{equation}
\label{eq:geo-normal}
\tan\varphi=\frac{Z_s/b^2}{p_s/a^2}
\qquad\Longrightarrow\qquad
Z_s=\frac{b^2}{a^2}\,p_s\tan\varphi=(1-e^2)\,p_s\tan\varphi .
\end{equation}
Substituting \eqref{eq:geo-normal} into \eqref{eq:geo-ellipse}, with $b^2=a^2(1-e^2)$,
\begin{equation}
\frac{p_s^2}{a^2}\Big(1+(1-e^2)\tan^2\varphi\Big)=1
\quad\Longrightarrow\quad
p_s^{2}=\frac{a^2\cos^2\varphi}{\cos^2\varphi+(1-e^2)\sin^2\varphi}
       =\frac{a^2\cos^2\varphi}{1-e^2\sin^2\varphi},
\end{equation}
that is,
\begin{equation}
\label{eq:geo-ps}
p_s=N(\varphi)\cos\varphi,
\qquad
N(\varphi)\equiv\frac{a}{\sqrt{1-e^2\sin^2\varphi}},
\end{equation}
and, back through \eqref{eq:geo-normal},
\begin{equation}
\label{eq:geo-zs}
Z_s=(1-e^2)\,N(\varphi)\sin\varphi .
\end{equation}
The quantity $N(\varphi)$ so defined is the \emph{prime-vertical radius of curvature}; a
short computation shows it is precisely the distance, along the normal, from the surface
point to the polar axis (the segment $Q$--$P$ of Figure~\ref{fig:ellipsoid}): the normal
direction in the meridian plane is $(\cos\varphi,\sin\varphi)$, and following it inward
from $(p_s,Z_s)$ a length $N$ lands on the axis,
$p_s-N\cos\varphi=0$ by \eqref{eq:geo-ps}.

\paragraph{Adding the height.}
The height $h$ of a fix is measured along the same normal, whose unit vector in ECEF
components is
\begin{equation}
\label{eq:geo-up}
\hat{\mathbf{u}}=(\cos\varphi\cos\lambda,\;\cos\varphi\sin\lambda,\;\sin\varphi).
\end{equation}
Adding $h\,\hat{\mathbf{u}}$ to the surface point
\eqref{eq:geo-ps}--\eqref{eq:geo-zs} in the meridian plane and restoring the longitude via
\eqref{eq:geo-meridian}:
\begin{equation}
X=(N+h)\cos\varphi\cos\lambda,\qquad
Y=(N+h)\cos\varphi\sin\lambda,\qquad
Z=\big(N(1-e^2)+h\big)\sin\varphi,
\end{equation}
which is Eq.~\eqref{eq:ecef}. The asymmetry between the horizontal factor $(N+h)$ and the
vertical one $\big(N(1-e^2)+h\big)$ is the ellipsoid's flattening at work: on a sphere
($e=0$) both reduce to $R+h$ and the geodetic and geocentric latitudes coincide.

\section{ECEF to local ENU}
\label{app:geodesy-enu}

The local frame at the origin fix $(\varphi_0,\lambda_0,h_0)$ is spanned by three unit
vectors: East, tangent to the parallel; North, tangent to the meridian, pointing to the
pole; Up, along the ellipsoid normal \eqref{eq:geo-up}. East is the direction in which the
position \eqref{eq:geo-meridian} changes under $\lambda$ alone,
\begin{equation}
\hat{\mathbf{e}}_E=\frac{\partial\mathbf{X}/\partial\lambda}
                        {\lVert\partial\mathbf{X}/\partial\lambda\rVert}
=(-\sin\lambda_0,\;\cos\lambda_0,\;0),
\end{equation}
Up is $\hat{\mathbf{u}}$ of \eqref{eq:geo-up} evaluated at
$(\varphi_0,\lambda_0)$, and North completes the right-handed triad,
\begin{equation}
\hat{\mathbf{e}}_N=\hat{\mathbf{e}}_U\times\hat{\mathbf{e}}_E
=(-\sin\varphi_0\cos\lambda_0,\;-\sin\varphi_0\sin\lambda_0,\;\cos\varphi_0).
\end{equation}
The ENU components of a displacement $\Delta\mathbf{X}$ are its projections on the triad,
$E=\hat{\mathbf{e}}_E\!\cdot\!\Delta\mathbf{X}$ and likewise for $N$ and $U$; stacking the
three row vectors gives exactly the matrix of Sec.~\ref{sec:enu},
\begin{equation}
\begin{pmatrix}E\\N\\U\end{pmatrix}=
\underbrace{\begin{pmatrix}
-\sin\lambda_0 & \cos\lambda_0 & 0\\
-\sin\varphi_0\cos\lambda_0 & -\sin\varphi_0\sin\lambda_0 & \cos\varphi_0\\
\cos\varphi_0\cos\lambda_0 & \cos\varphi_0\sin\lambda_0 & \sin\varphi_0
\end{pmatrix}}_{\textstyle R\,}
\Delta\mathbf{X},
\end{equation}
whose rows are orthonormal by construction, so $R$ is a rotation
($R^{\mathsf T}R=\mathbb{1}$)~\cite{esa_navipedia_ellipsoidal_cartesian}. Two remarks used
in the body follow directly. First, the rows depend on the \emph{geodetic} latitude
because Up is the ellipsoid normal --- the same $\varphi$ that enters
Eq.~\eqref{eq:ecef}, which is why no latitude conversion appears anywhere in the
pipeline. Second, the matrix is evaluated once per flight, at the origin fix: the frame
is a fixed tangent frame, not one that travels with the wing.

\section{Size of the tangent-plane error}
\label{app:geodesy-error}

Working in a fixed tangent frame replaces distances measured \emph{along} the curved
surface by straight-line (chord) distances in the plane. For two surface points separated
by an arc $d$ on a sphere of radius $R_\oplus$, the chord is
\begin{equation}
c=2R_\oplus\sin\!\frac{d}{2R_\oplus}
 = d-\frac{d^{3}}{24\,R_\oplus^{2}}+O\!\big(d^{5}\big),
\end{equation}
so the horizontal distance error of the tangent-plane picture is
$d^{3}/(24R_\oplus^{2})$: about \SI{1}{\centi\metre} at $d=\SI{20}{\kilo\metre}$ and
\SI{13}{\centi\metre} at $d=\SI{50}{\kilo\metre}$ --- far below the metre-scale GPS noise
on the same displacement, which is the comparison made in Sec.~\ref{sec:enu}. The
\emph{vertical} coordinate of a distant surface point is a different matter: the surface
falls below the tangent plane by the sagitta $\approx d^{2}/(2R_\oplus)$, which reaches
tens of metres already at $d=\SI{20}{\kilo\metre}$. This is why the analysis never takes
its working altitude from the plane geometry: the vertical coordinate is the barometric
altitude itself (Sec.~\ref{sec:enu}), and the $U$ component serves only the geometric
construction.
\end{revblock}
